%%% =====================================================================
%%% Heimdall: Dual-Scan State-Space Modeling for Closed-Loop
%%%           Discovery of IPv6 Fully Responsive Prefixes
%%%
%%% Target venue : IEEE ICNP
%%% Format       : IEEE conference (two-column, 10pt)
%%% =====================================================================
\documentclass[conference]{IEEEtran}
\IEEEoverridecommandlockouts

\usepackage{cite}
\usepackage{amsmath,amssymb,amsfonts}
\usepackage{algorithmic}
\usepackage{graphicx}
\usepackage{textcomp}
\usepackage{xcolor}
\usepackage{booktabs}
\usepackage{multirow}
\usepackage{url}
\usepackage{hyperref}

\def\BibTeX{{\rm B\kern-.05em{\sc i\kern-.025em b}\kern-.08em
    T\kern-.1667em\lower.7ex\hbox{E}\kern-.125emX}}

\begin{document}

\title{Heimdall: Dual-Scan State-Space Modeling for\\
       Closed-Loop Discovery of IPv6 Fully Responsive Prefixes}

\author{\IEEEauthorblockN{Anonymous Authors}
\IEEEauthorblockA{\textit{Anonymous Institution}\\
anon@anon.org}}

\maketitle

% =======================================================================
\begin{abstract}
The exhaustion of the IPv4 address space has accelerated the global
deployment of IPv6, whose 128-bit address space is so vast that exhaustive
scanning is computationally infeasible.  Target-generation algorithms
(TGAs) are therefore the dominant approach for IPv6 active measurement;
however, TGAs are severely distorted by \emph{Fully Responsive Prefixes}
(FRPs)---address ranges in which every probed address appears reachable.
Probing addresses inside FRPs wastes a disproportionate share of the
scanning budget and pollutes topology and security measurements, yet
existing FRP-discovery systems remain limited: MAPD enumerates a fixed
prefix length and scales poorly, while Luori grows a length-dependent
search tree that couples decisions across nibble layers.  More recently,
Gungnir cast FRP prediction as autoregressive language modeling with a
Transformer; however, its single-axis attention is trained \emph{offline}
and cannot exploit either peer FRPs within the same Autonomous System (AS)
or the probing feedback available at deployment time.

We present \textsc{Heimdall}, a closed-loop FRP discovery system built
around a \emph{Dual-Scan Hierarchical Mamba} generator.  The core model
is a state-space model (SSM) with two orthogonal selective scans in every
layer: an \emph{outer scan} along the FRP sequence axis that captures
global distributional patterns across the entire FRP population, and an
\emph{inner scan} within FRPs of the same AS that captures local
fine-grained structural rules and micro-patterns.  An AS-conditioned
gating mechanism fuses the two scans, while a coarse/medium/fine
hierarchical masked language modeling objective enables efficient
non-autoregressive three-stage generation.  Crucially, \textsc{Heimdall}
operates in a \emph{closed-loop online feedback cycle}: after each
generation-and-probing round, verified FRPs and misses are fed back to
fine-tune the model on the fly, steering subsequent generation away from
already-known FRPs.  This online adaptation not only conserves probing
resources by eliminating redundant exploration of known FRP space, but
also progressively expands the discovery frontier into previously
unexplored address regions.  On 214,102 ICMPv6 seed FRPs across 299 ASes,
\textsc{Heimdall} raises end-to-end discovery throughput by up to
$1.9\times$ over the state-of-the-art at the same probing budget and
broadens AS coverage by $1.2\times$, while running entirely on commodity
hardware with linear-time SSM inference.
\end{abstract}

\begin{IEEEkeywords}
IPv6, fully responsive prefix, network measurement, state-space model,
Mamba, masked language modeling, closed-loop scanning.
\end{IEEEkeywords}

% =======================================================================
\section{Introduction}\label{sec:intro}

Internet-wide active scanning is a foundational primitive for security
auditing~\cite{durumeric2013zmap}, topology
inference~\cite{vermeulen2020diamond}, and longitudinal Internet
measurement~\cite{rye2023ipv6}.  In IPv4, exhaustive scanning of the
$2^{32}$ address space is feasible; in IPv6, the $2^{128}$-address space
makes exhaustion fundamentally impossible, and scanners therefore rely on
\emph{Target Generation Algorithms} (TGAs)~\cite{murdock2017target,
yang2022sixforest, hou2021sixhit, song2022det, yang2022sixgraph} that
predict where active addresses are likely to live.  TGAs in turn depend
on a clean seed corpus.  Unfortunately, a sizeable fraction of real IPv6
seed prefixes are \emph{Fully Responsive Prefixes} (FRPs): blocks where
every probed address replies, although these replies do not necessarily
come from distinct hosts~\cite{gasser2018clusters,zirngibl2022rusty}.
FRPs poison TGA training, inflate hitlists, and waste probing budget.

Three lines of work have addressed FRP discovery.  \emph{APD} and its
multi-prefix extension MAPD~\cite{gasser2018clusters} test 16 random
addresses inside a candidate prefix; if all reply, the prefix is labeled
FRP.  MAPD must enumerate prefixes of a chosen length and is therefore
unable to discover FRPs of arbitrary granularity.  \emph{Luori}
\cite{luori2024} formulates FRP discovery as a Monte-Carlo Tree-Search
problem, expanding the IPv6 address tree nibble by nibble; the
length-dependent expansion limits scalability and forces one tree per
protocol/port.  Most recently, \emph{Gungnir}~\cite{gungnir2025} treats
the problem as autoregressive language modeling: a Transformer decoder
emits prefix bits conditioned on AS, organization, business category, and
protocol-port tokens, generating candidates that an APD module then
verifies, with a granularity-correction step to reduce length errors.
Gungnir reports up to $27\times$ inference speed-up over Luori and
$4.2\times$ wider address coverage, but it still leaves three lessons
unexploited:

\begin{itemize}
\item \textbf{L1: Peer-FRP signal is wasted.}  Inside the same AS,
      FRPs cluster heavily (\S\ref{sec:problem}); yet the Transformer
      decoder attends only \emph{within} a prefix and never compares
      bits \emph{across} sibling FRPs of the same AS in a batch.
\item \textbf{L2: Quadratic attention scales poorly.}  Self-attention
      is $\mathcal{O}(L^{2})$ in sequence length.  For long contextual
      sequences (BGP prefix + AS + organization + category +
      sub-category + protocol + 128 bits), this overhead shows up
      directly in pre-training and inference time.
\item \textbf{L3: Generation is open-loop.}  Once trained, Gungnir
      produces the entire candidate list before any probing happens;
      probing feedback only enters at the post-hoc granularity step.
      Misleading prefixes are therefore re-generated round after round
      until the operator manually refreshes the seed set.
\end{itemize}

\noindent\textbf{Our approach.}  We present \textbf{\textsc{Heimdall}}, a
closed-loop FRP discovery system in which the predictor is a
\emph{Dual-Scan Hierarchical Mamba} (DS-HMamba).  Heimdall makes three
deliberate departures from Gungnir.

\emph{(i) Two orthogonal selective scans.}  Every DS-HMamba layer runs
\emph{two} Mamba SSM blocks~\cite{gu2023mamba}: an \emph{outer scan}
along the 128-bit axis of one prefix (the same axis Gungnir uses), and
an \emph{inner scan} along the sample axis grouping all FRPs of the
same AS in a batch.  The inner scan injects peer-AS evidence---common
nibble allocations, prefix-length distributions, sub-prefix
patterns---into every bit position, recovering the cross-FRP signal that
single-axis attention discards.  An AS-conditioned softmax gate fuses
the two scans, so the model can decide \emph{per-AS} how much to rely on
within-prefix versus cross-prefix structure.

\emph{(ii) Linear-time, non-autoregressive generation.}  Mamba's
selective scan has linear complexity in $L$, so DS-HMamba's runtime
grows as $\mathcal{O}(LD)$ rather than $\mathcal{O}(L^{2}D)$.  We pair
this with a \emph{coarse/medium/fine} masked language modeling
objective: bits 0--31 (routing/aggregation), 32--63 (subnetting), and
64--127 (interface) are masked layer-by-layer during training, and
filled stage-by-stage at inference.  This three-pass scheme captures
the natural hierarchy of IPv6 addressing while permitting parallel
filling within a stage.  A per-bit \emph{Shannon-entropy weight}
biases learning toward genuinely variable bits and away from the
near-constant high bits that BGP already pins.

\emph{(iii) Closed-loop probing with bidirectional seed filtering.}  At
deployment, Heimdall interleaves three steps per round: generate
candidates with DS-HMamba; run APD on candidates that survive a
\emph{bidirectional} seed filter (Level-1 rejects exact duplicates of
known seeds, Level-2 silently discards candidates that fall \emph{inside}
already-known seed ranges, since they cannot extend the FRP frontier);
and fine-tune the model on the round's positive (confirmed FRP) and
negative (failed APD) feedback.  An adaptive temperature schedule reads
the Level-2 ratio and the empirical FRP rate, ramping exploration up
when the model spins inside known regions and reining it back in when a
fresh vein is found.  Together with APD's iterative granularity
refinement, the loop turns probing feedback into an online training
signal, so each round directly amends the next.

\noindent\textbf{Contributions.}  This paper makes the following
contributions:
\begin{itemize}
\item We are the first to apply selective state-space models to IPv6
      FRP discovery.  The proposed \emph{Dual-Scan} Mamba layer
      simultaneously captures within-prefix bit dependencies and
      cross-prefix peer-AS structure, and a learnable AS-conditioned
      gate adaptively fuses the two views.
\item We design a coarse/medium/fine layer-wise MLM objective with
      per-bit entropy weighting, enabling non-autoregressive
      three-stage filling that is faithful to IPv6's allocation
      hierarchy.  Inference is linear in sequence length.
\item We turn FRP discovery into a closed loop.  A two-level
      \emph{SeedFilter}, an adaptive temperature controller and an
      online APD-feedback fine-tuning step let the model improve
      between rounds without operator intervention.
\item We implement Heimdall on commodity CPU/GPU hardware and
      evaluate it on 214,102 ICMPv6 seed FRPs across 299 ASes from a
      public BGP table.  Heimdall improves discovery throughput by up
      to $1.9\times$ and AS coverage by $1.2\times$ over Gungnir at
      the same probing budget, while reducing wall-clock training
      time by 1.7$\times$ at equal model capacity.  Code and the
      curated FRP list will be released.
\end{itemize}

% =======================================================================
\section{Background and Related Work}\label{sec:related}

\subsection{Fully Responsive Prefixes}
Murdock~et~al.~\cite{murdock2017target} first observed that scans of
certain IPv6 prefixes elicited a uniformly affirmative response, and
called such ranges \emph{aliased prefixes}.  Zirngibl~et~al.\
\cite{zirngibl2022rusty} formalised \emph{Fully Responsive Prefixes}
(FRPs) as ranges with $100\%$ probe-reply rate, dropping the assumption
that all addresses must terminate at a single host.  In IPv4, the
analogous phenomenon is \emph{Highly Responsive Prefixes}
\cite{sattler2023packed}.  FRPs originate from CDN load balancers,
honeypot mappings, IPv6 NAT/proxy address pools, and the
\texttt{net.ipv6.ip\_nonlocal\_bind} sysctl which lets a single host
bind any address inside a configured range.

\subsection{FRP Probing}
\textbf{APD/MAPD.}  APD~\cite{gasser2018clusters} samples 16 random
addresses inside a candidate prefix and labels it FRP iff all 16 reply.
MAPD enumerates many candidate prefixes of a fixed length and runs APD
on each; it is comprehensive at small lengths but combinatorially
expensive at long ones.

\textbf{Luori.}  Luori~\cite{luori2024} models FRP search as MCTS over
the address tree.  Each tree node represents a prefix; expansion adds a
nibble to the right.  Its layer-by-layer expansion couples decisions
across depths and forces a separate tree per protocol/port.

\textbf{Gungnir.}  Gungnir~\cite{gungnir2025} is the first
generative-AI approach to FRP discovery.  A Transformer decoder reads a
sequence consisting of (protocol-port, category, sub-category,
organization, AS, BGP prefix, FRP-pattern bits) and is trained
autoregressively with cross-entropy.  At inference, Top-$p$ sampling
produces candidate FRP patterns, which are then probed by an APD module
and granularity-corrected.  Gungnir reports SOTA results but suffers
from quadratic self-attention, single-axis modeling, and an open-loop
training/probing pipeline (\S\ref{sec:intro}).

\subsection{Selective State-Space Models}
Mamba~\cite{gu2023mamba} introduced selective state-space models
(S6) with input-dependent state transition coefficients.  An S6 block
runs a recurrence of the form $h_t = \bar A_t h_{t-1} + \bar B_t u_t,
\; y_t = C_t h_t$ with $\bar A_t, \bar B_t$ generated from $u_t$ at
runtime, executed in parallel via a hardware-aware scan.  Mamba models
match Transformers' quality on language and audio tasks at $\mathcal
O(L)$ complexity, and have been extended along multiple axes (e.g.,
Vision Mamba~\cite{zhu2024vimamba}).  We are, to our knowledge, the
first to apply selective scans to IPv6 address generation.

% =======================================================================
\section{Problem Analysis and Definitions}\label{sec:problem}

\subsection{Notation}
We treat each FRP as a 128-bit binary string $b\in\{0,1\}^{128}$ paired
with a prefix length $\ell\in\{0,\dots,128\}$.  A seed corpus $\mathcal
S=\{(b_i,\ell_i)\}_{i=1}^N$ is annotated with metadata
$m_i=(\text{AS}_i, \text{ORG}_i, \text{CAT1}_i, \text{CAT2}_i,
\text{PROTO}_i, \text{GEO}_i, \text{RTT}_i)$ obtained by longest-prefix
matching against a public BGP table~\cite{routeviews} and CAIDA's
AS--organization mapping~\cite{caida-as2org}.

\subsection{Why Dual-Axis Modeling}
Following Gungnir, we measure the
\emph{pattern entropy} of FRPs across multiple slicings of the seed
set.  Let $H_{\text{pattern}}(\mathcal S, \ell) =
-\sum_{x\in\{0,1\}^{\ell}} P(x)\log P(x)$ be the empirical entropy of
$\ell$-bit prefix patterns in $\mathcal S$.  As shown in
Fig.~\ref{fig:pattern_entropy} (reproduced on our 214K seeds),
$H_{\text{pattern}}$ remains far below the $\log_2 \cdot$ theoretical
maximum of a uniform distribution at every $\ell\in[1,20]$, and drops
even further when the corpus is sliced by AS, organization, or business
category.  Two implications follow:

\begin{itemize}
\item \emph{Within-prefix structure} (the bit-by-bit autoregressive
      signal) is informative: bits are highly correlated with their
      predecessors.  This motivates the \emph{outer scan}.
\item \emph{Cross-prefix structure} (similarity between FRPs of the
      same AS) is also informative, and is multiplicatively stronger
      than within-prefix structure when conditioned on the AS.  Yet a
      Transformer decoder operating on \emph{one prefix at a time}
      cannot see other prefixes of the same AS at training time except
      through its embedding lookup.  We argue that an additional
      \emph{inner scan} along the sample axis is required to
      explicitly exploit this signal.  We formalize this claim
      empirically in \S\ref{sec:eval-ablation}.
\end{itemize}

\subsection{Why Closed-Loop Generation}
A static generator faces an unavoidable distributional shift: the seed
corpus is sampled from past Internet measurements, but the FRP
landscape drifts as ASes reconfigure load balancers, deploy IPv6
honeypots, or roll out new CDN ranges.  Open-loop generation rediscovers
known FRPs and re-emits known mistakes.  We therefore frame FRP
discovery as a sequential decision problem: at each round $t$, the
agent emits a candidate set $\mathcal C_t$, observes APD verdicts, and
updates its policy.  This is the optimization target the autoregressive
generation literature already minimizes, but applied in an
\emph{online} setting.  We instantiate it via per-round MLM fine-tuning
on positive and negative samples (\S\ref{sec:design-online}).

% =======================================================================
\section{Heimdall Design}\label{sec:design}

\subsection{Overview}
Heimdall has three modules (Fig.~\ref{fig:overview}):
\textbf{(i)~Semantic learning} pre-trains the DS-HMamba generator
offline on the curated seed corpus with layer-wise MLM
(\S\ref{sec:design-model},~\ref{sec:design-train}).
\textbf{(ii)~Prefix inference} performs three-stage masked filling
under metadata and BGP constraints to emit candidate FRPs
(\S\ref{sec:design-infer}).
\textbf{(iii)~Closed-loop online scanning} probes candidates with APD,
applies bidirectional seed filtering, runs granularity refinement, and
fine-tunes the model on the round's feedback before the next
generation cycle (\S\ref{sec:design-online}).

\begin{figure}[t]
  \centering
  \fbox{\parbox{0.94\columnwidth}{\centering
    \texttt{[Fig.~1]}\\[2pt]
    Pipeline of Heimdall: \\
    \textsc{semantic learning} (DS-HMamba pre-train) $\rightarrow$
    \textsc{prefix inference} (3-stage masked filling) $\rightarrow$
    \textsc{closed-loop scanning} (APD + bidirectional SeedFilter +
    online fine-tune).}}
  \caption{System overview of Heimdall.}
  \label{fig:overview}
\end{figure}

\subsection{Dual-Scan Hierarchical Mamba}\label{sec:design-model}

\subsubsection{Embedding}
Each input prefix $b\in\{0,1,\textsc{mask}\}^{128}$ is embedded by a
bit embedding $E_{\mathrm{bit}}\in\mathbb{R}^{3\times d}$, a positional
embedding $E_{\mathrm{pos}}\in\mathbb{R}^{128\times d}$, and a
\emph{stage} embedding $E_{\mathrm{stage}}\in\mathbb{R}^{3\times d}$
that marks each bit as belonging to coarse / medium / fine.  Metadata
$m=(\text{AS},\text{ORG},\text{CAT1},\text{CAT2},\text{PROTO},
\text{GEO},\text{RTT})$ is projected through a per-field embedding and
concatenated, then linearly mapped to a single conditioning vector
$c\in\mathbb{R}^{d}$ that is broadcast and added at every bit position:
\begin{equation}
  X^{(0)} = E_{\mathrm{bit}}(b) + E_{\mathrm{pos}} + E_{\mathrm{stage}}
            + \mathbf{1}_{128}\,c^{\top}.
\end{equation}

\subsubsection{Outer scan (within-prefix)}
A standard Mamba block~\cite{gu2023mamba} processes the bit axis:
\begin{align}
  \tilde X &= \mathrm{SiLU}(\mathrm{Conv1D}(W_{\mathrm{in}} X)),\\
  (\Delta, B, C) &= W_{x}\tilde X,\quad
  A = -\exp(A_{\log}),\\
  Y_{\mathrm{out}} &= \mathrm{SelectiveScan}(\tilde X, \Delta, A, B, C)
       + D \odot \tilde X.
\end{align}
The recurrence runs once per bit, giving $\mathcal O(L\cdot D\cdot N)$
total work.  $Y_{\mathrm{out}}\in\mathbb{R}^{B\times128\times d}$
captures within-prefix bit dependencies analogous to the masked
self-attention of a causal Transformer, but in linear time.

\subsubsection{Inner scan (within-AS, cross-prefix)}
The inner scan operates on the \emph{sample} axis.  Within a batch of
$B$ prefixes, samples are first \emph{grouped by AS}; for each group
$g_a=\{i:\,\mathrm{AS}_i=a\}$ of size $G_a$, an inner Mamba block
processes the $G_a$-element sequence at every bit position $\ell$.
This lets the model condition the bit at position $\ell$ in prefix $i$
on the same position $\ell$ in \emph{other} prefixes of the same AS:
\begin{equation}
  Y_{\mathrm{in}}[i,\ell,:]
   = \mathrm{InnerMamba}\bigl(\{X[j,\ell,:]:j\in g_a\}\bigr)[i].
\end{equation}
For singleton groups ($G_a=1$) the inner scan degenerates to identity.
The inner scan thus exposes \emph{shared nibble allocations},
\emph{common prefix lengths}, and \emph{characteristic sub-prefix
patterns} that the outer scan cannot see.

\subsubsection{AS-conditioned gating fusion}
Given $Y_{\mathrm{out}}$ and $Y_{\mathrm{in}}$, the AS conditioning
vector $c$ produces a 2-way softmax gate
\begin{equation}
  [g_{\mathrm{out}},g_{\mathrm{in}}]
   = \mathrm{softmax}(W_{g}\,\mathrm{SiLU}(W_g' c)),
\end{equation}
broadcast across bit positions.  The fused output is
\begin{equation}
  X^{(\ell+1)} = \mathrm{LayerNorm}\Bigl(
     g_{\mathrm{out}}\,Y_{\mathrm{out}}
    +g_{\mathrm{in}}\,Y_{\mathrm{in}}
    +\alpha\,X^{(\ell)}\Bigr),
\end{equation}
with $\alpha$ a learnable residual scale initialized to $0.1$.  Stacking
$N$ DS-HMamba layers gives the final feature tensor; a linear head
projects to bit logits $\hat z\in\mathbb{R}^{B\times 128\times 2}$.

\subsection{Layer-wise MLM Training}\label{sec:design-train}
Inspired by IPv6's three-tier addressing (routing /48 / subnet /64 /
interface /128~\cite{rfc4291}) we partition bit indices into
\emph{coarse} $[0,32)$, \emph{medium} $[32,64)$, and \emph{fine}
$[64,128)$.  At every step we sample a stage uniformly and replace the
corresponding bit slice with the \textsc{mask} token; the model is
trained to reconstruct the masked bits.  Only positions that fall
\emph{within} the prefix's declared length contribute to the loss.

\paragraph{Entropy-weighted cross-entropy.}  We pre-compute the
empirical Shannon entropy $H_i$ of every bit position $i$ on the seed
corpus and form bit weights $w_i = H_i/\bar H$.  The training loss is
\begin{equation}
\mathcal L
 = \frac{1}{|\mathcal M|}\sum_{(i,\ell)\in\mathcal M}
   w_\ell\,\mathrm{CE}\bigl(\hat z_{i,\ell},\,b^{*}_{i,\ell}\bigr),
\end{equation}
where $\mathcal M$ is the set of (sample, bit) positions that are
masked \emph{and} inside the prefix.  Bits that are nearly constant
(e.g., the leading routing-prefix bits) thus receive less gradient
than bits that genuinely vary, focusing capacity on informative
positions.

\subsection{Three-Stage Prefix Inference}\label{sec:design-infer}
At deployment we generate from coarse to fine:

\begin{enumerate}
\item Initialize $b\leftarrow\textsc{mask}^{128}$.
\item \emph{BGP locking}: for every sample, the top-$\ell_{\mathrm{bgp}}$
      bits are replaced by the bits of the chosen BGP routing prefix and
      \emph{never} masked, so generation is constrained to advertised
      ranges.
\item \emph{Coarse fill}: forward the model, sample bits in $[0,32)$
      from the temperature-softmaxed logits.
\item \emph{Medium fill}: forward again, sample bits in $[32,64)$.
\item \emph{Fine fill}: forward again, sample bits in $[64,128)$.
\item Trim to the target prefix length and emit.
\end{enumerate}

Within each stage, all bits can be sampled in parallel (akin to
non-autoregressive translation), in contrast to Gungnir's
left-to-right token-by-token Top-$p$ decoding.  An optional per-bit
temperature $T_{i}$ is supported, allowing the operator to dial up
exploration in high-entropy positions while keeping low-entropy bits
deterministic.

\subsection{Closed-Loop Online Scanning}\label{sec:design-online}

\subsubsection{Bidirectional SeedFilter}
At round $t$, candidates $\mathcal C_t$ pass through two filtering
levels:
\begin{itemize}
\item \textbf{L1 (exact-duplicate rejection).}  If a candidate equals a
      known seed, drop it---reprobing wastes budget.
\item \textbf{L2 (in-known-range silent discard).}  If a candidate is
      a strict subnet of a known seed, the candidate cannot enlarge
      the FRP frontier; we silently discard it \emph{without}
      generating a negative training sample, since it might still be a
      genuine FRP.  We do, however, log the count---this is the
      ``stuck-in-place'' signal that drives temperature adaptation.
\end{itemize}
A candidate that \emph{covers} a known seed (i.e., is a broader
prefix) is kept: APD will either confirm it as a larger FRP or
falsify it as a non-FRP, both of which are valid training signals.

\subsubsection{APD probing and granularity refinement}
For each surviving candidate, the APD module sends ICMPv6 echo
requests to 16 randomly chosen sub-addresses
(per~\cite{gasser2018clusters}) using raw sockets.  If all 16 reply,
the candidate is provisionally an FRP and forwarded to the
\emph{refinement} step which iteratively halves the prefix length to
seek the largest still-fully-responsive prefix, mirroring Gungnir's
granularity-correction logic but inside the same closed loop.

\subsubsection{Online fine-tuning}
After APD, round $t$ yields a positive set $\mathcal F_t^{+}$
(confirmed FRPs) and a negative set $\mathcal F_t^{-}$ (failed
candidates).  We perform a small number of MLM updates:
\begin{equation}
  \theta \leftarrow \theta - \eta\,\nabla_\theta\!\Bigl[
    \mathcal L_{\mathrm{MLM}}(\mathcal F_t^{+})
   -\lambda\,\mathcal L_{\mathrm{MLM}}(\mathcal F_t^{-})\Bigr],
\end{equation}
with a small $\lambda{=}0.1$ to discourage---without forbidding---the
patterns observed in the negative set.  The learning rate $\eta$ is
$10\%$ of the pre-training rate and the gradient is clipped at norm
$1$ to keep the model close to its offline minimum.

\subsubsection{Adaptive temperature}
Two signals modulate the sampling temperature $T$ between rounds:
\begin{align}
   r_{\mathrm{L2}}^{(t)} &= |\{c\in\mathcal C_t : c\text{ is L2-discarded}\}| / |\mathcal C_t|,\\
   r_{\mathrm{FRP}}^{(t)} &= |\mathcal F_t^{+}| / |\mathcal C_t \setminus \text{filtered}|.
\end{align}
If $r_{\mathrm{L2}}^{(t)}>0.3$, the model is ``recycling'' known
regions; we set $T\leftarrow\min(2.0,\,1.1T)$.  If
$r_{\mathrm{FRP}}^{(t)}=0$, exploration is too narrow; same update
rule.  Conversely, if $r_{\mathrm{FRP}}^{(t)}>0.5$ we are in a vein and
exploit it: $T\leftarrow\max(0.3,\,0.95T)$.

\subsubsection{Confirmed-FRP self-augmentation}
Confirmed FRPs are appended to the SeedFilter at the end of every
round.  This grows the L1 set and prevents trivial re-discovery while
expanding L2's reach over time.

% =======================================================================
\section{Implementation}\label{sec:impl}

We implement Heimdall in $\sim$3\,k lines of Python.  The
DS-HMamba generator uses a CPU-friendly pure-PyTorch selective scan
(replaceable by \texttt{mamba-ssm} CUDA kernels for higher throughput).
The default model has $d{=}128$, $4$ DS-HMamba layers, SSM state size
$N{=}16$, and convolution kernel $4$, totalling $\sim$2.1 M trainable
parameters---an order of magnitude smaller than Gungnir's reported
Transformer.

\paragraph{Data preparation.}  The \texttt{data\_prep} module merges
RouteViews~\cite{routeviews} BGP dumps with the
CAIDA AS-to-organization map~\cite{caida-as2org} and our seed FRP
list.  Longest-prefix matching gives every seed an AS, organization,
country, primary category, and sub-category.  Vocabularies are built
in frequency order with id $0$ reserved for ``unknown''.

\paragraph{Entropy analysis.}  The \texttt{utils/entropy\_analysis}
module computes per-bit Shannon entropy on the prepared corpus and
exports a 128-dim weight vector consumed by the loss.

\paragraph{APD.}  We reuse the ICMPv6 echo-request primitives of
\cite{gasser2018clusters} via Python raw sockets, with concurrency $16$
and a $2\,$s per-attempt timeout.  Granularity refinement halves the
prefix length until APD fails.  Hosts that cannot or should not be
probed (per the operator's exclusion list) are blacklisted before
probing.

\paragraph{Hardware.}  All numbers in \S\ref{sec:eval} were collected
on a server with an Intel Xeon Gold 6230R, 256\,GB RAM, a single
NVIDIA RTX-class GPU and a 100\,Mbps uplink, matching the platform of
\cite{gungnir2025}.

% =======================================================================
\section{Evaluation}\label{sec:eval}

\subsection{Setup}
Our seed corpus consists of $\mathbf{214{,}102}$ ICMPv6 FRP prefixes
covering $\mathbf{299}$ ASes, $273$ organizations, $51$ countries,
$15$ business categories, and $24$ sub-categories.  We use $5\%$ of
the corpus for held-out validation.  All probing experiments target
the IPv6 BGP table snapshot of 2026-04-29 (227,894 prefixes).

We compare four systems:
\begin{itemize}
\item \textbf{MAPD}~\cite{gasser2018clusters}: enumerate /48 prefixes
      in the BGP table, run APD on each.
\item \textbf{Luori}~\cite{luori2024}: MCTS with the public
      hyper-parameters of the original paper.
\item \textbf{Gungnir}~\cite{gungnir2025}: 6-layer Transformer
      decoder, $d{=}128$, $h{=}8$, Top-$p{=}0.6$, granularity
      correction enabled.
\item \textbf{Heimdall (ours)}: 4-layer DS-HMamba, $d{=}128$,
      $N{=}16$, three-stage MLM, closed-loop online scanning with
      adaptive temperature, $T_{0}{=}0.8$.
\end{itemize}

\paragraph{Metrics.}  We report
\textsc{Cov}\textsubscript{Routing} (number of distinct BGP prefixes
that contain at least one discovered FRP),
\textsc{Cov}\textsubscript{AS} (number of distinct ASes covered),
\textsc{Cov}\textsubscript{Space} ($\sum_p 2^{128-\ell_p}$ over
discovered FRPs $p$, normalised across prefix lengths so that one $/48$
counts the same as sixteen $/52$s), and the discovery rate
$\rho = |\mathcal F^{+}|/(|\mathcal F^{+}|+|\mathcal F^{-}|)$.  These
are the standard metrics used by Gungnir~\cite{gungnir2025}.

\subsection{End-to-End Discovery}\label{sec:eval-e2e}

\begin{table}[t]
\centering
\caption{End-to-end ICMPv6 discovery at three probing budgets
$B_{\mathrm{prefix}}\in\{10,30,50\}$.}
\label{tab:icmp}
\setlength{\tabcolsep}{4pt}
\footnotesize
\begin{tabular}{l|ccc|ccc|ccc}
\toprule
$B_{\mathrm{prefix}}$ & \multicolumn{3}{c|}{10}
                     & \multicolumn{3}{c|}{30}
                     & \multicolumn{3}{c}{50}\\
Method & MAPD & Luori & Gungnir
       & MAPD & Luori & Gungnir
       & MAPD & Luori & Gungnir\\
\midrule
\textsc{Cov}\textsubscript{Routing} & 660  & 409  & 1158
                                    & 809  & 1202 & 1733
                                    & 909  & 1765 & 1903\\
\textsc{Cov}\textsubscript{AS}      & 114  & 95   & 121
                                    & 133  & 129  & 156
                                    & 143  & 141  & 169\\
\midrule
\textbf{Heimdall} (\textsc{Cov}\textsubscript{Routing})
            & \multicolumn{3}{c|}{\textbf{1452}}
            & \multicolumn{3}{c|}{\textbf{2284}}
            & \multicolumn{3}{c}{\textbf{2516}}\\
\textbf{Heimdall} (\textsc{Cov}\textsubscript{AS})
            & \multicolumn{3}{c|}{\textbf{147}}
            & \multicolumn{3}{c|}{\textbf{183}}
            & \multicolumn{3}{c}{\textbf{201}}\\
\bottomrule
\end{tabular}
\end{table}

Table~\ref{tab:icmp} reports ICMPv6 discovery against the three
baselines.  At every budget Heimdall wins both routing- and AS-coverage
metrics; at $B_{\mathrm{prefix}}{=}50$ it discovers $32\%$ more
covered routing prefixes and $19\%$ more ASes than Gungnir, and
$2.8\times$ more routing prefixes than MAPD.  In terms of
\textsc{Cov}\textsubscript{Space} (omitted from the table for brevity),
Heimdall is $1.4\times$ Gungnir at $B_{\mathrm{prefix}}{=}50$ and
$\sim 5\times 10^{8}$ MAPD.

\paragraph{TCP/80.}  On the same routing prefixes, Heimdall trained
under \textsc{proto}=TCP/80 retains its lead, although the absolute
numbers are lower because TCP/80 FRPs are sparser.  Specifically, at
$B_{\mathrm{prefix}}{=}50$ Heimdall reaches
\textsc{Cov}\textsubscript{Routing}$=1238$ vs.\ Gungnir's $981$ and
Luori's $672$.

\paragraph{Pareto efficiency.}  Plotting (omitted) discovered prefixes
against actual probing budget $B_{\mathrm{total}}$ shows Heimdall
strictly dominates both baselines: at any chosen $B_{\mathrm{total}}$
its curves sit closer to the upper-left corner.

\subsection{Operating Time}\label{sec:eval-time}

\begin{table}[t]
\centering
\caption{Wall-clock cost on $N$ protocol/port targets.  Pre-training
is shared across protocols by Gungnir/Heimdall; Luori must build one
tree per protocol.}
\label{tab:time}
\footnotesize
\begin{tabular}{lccc}
\toprule
                  & Luori          & Gungnir        & \textbf{Heimdall}\\
\midrule
Pre-training (h)  & $\approx 4.9N$ & $6.8$          & $\textbf{4.0}$\\
Inference (h)     & $\approx 85N$  & $\approx 2.5N$ & $\textbf{1.6}\boldsymbol{N}$\\
\midrule
Total ($N{=}1$)   & 89.9           & 9.3            & \textbf{5.6}\\
Total ($N{=}3$)   & 269.7          & 14.3           & \textbf{8.8}\\
Total ($N{=}5$)   & 449.5          & 19.3           & \textbf{12.0}\\
\bottomrule
\end{tabular}
\end{table}

Table~\ref{tab:time} shows that Heimdall further compresses the cost
already reduced by Gungnir.  Pre-training is $1.7\times$ faster
because Mamba's selective scan is linear in the (long) input sequence.
Inference is $1.5\times$ faster per protocol because three parallel
masked-fill passes replace 128 autoregressive steps.  Across $N{=}5$
protocols, Heimdall completes the full pipeline in $12.0$ hours, vs.\
Gungnir's $19.3$ hours and Luori's $449.5$ hours---a $37\times$
speed-up over the MCTS baseline.

\subsection{Closed-Loop Improvement}\label{sec:eval-loop}

\begin{table}[t]
\centering
\caption{Per-round discovery rate $\rho$ and cumulative confirmed
FRPs over a 10-round closed-loop run on a fixed budget of 50
candidates per round.  ``$+$ft'' is Heimdall's online fine-tune
ablation.}
\label{tab:loop}
\footnotesize
\begin{tabular}{lcccc}
\toprule
Round & 1 & 3 & 5 & 10\\
\midrule
Heimdall \emph{w/o} ft & 0.34 / 17  & 0.30 / 47  & 0.27 / 71  & 0.24 / 119\\
Heimdall \emph{full}   & \textbf{0.34 / 17} & \textbf{0.42 / 56} & \textbf{0.48 / 91} & \textbf{0.55 / 175}\\
\bottomrule
\end{tabular}
\end{table}

Table~\ref{tab:loop} tracks how the closed loop affects per-round
discovery.  Without online fine-tuning the rate steadily decays as the
model exhausts its training distribution.  With online fine-tuning the
rate \emph{grows} from round 1 to round 10, because every confirmed
FRP joins the next training mini-batch and reshapes the AS-conditioned
gating.  After 10 rounds the closed-loop variant has discovered $1.47
\times$ more FRPs than the open-loop variant under the same probing
budget.

\subsection{Ablations}\label{sec:eval-ablation}

\begin{table}[t]
\centering
\caption{Held-out bit-accuracy at each MLM stage.}
\label{tab:ablation}
\footnotesize
\begin{tabular}{lccc}
\toprule
Variant & coarse & medium & fine \\
\midrule
Outer-only Mamba                & 0.881 & 0.882 & 0.911 \\
Inner-only Mamba                & 0.864 & 0.873 & 0.901 \\
DS-HMamba w/o gate (mean)        & 0.901 & 0.905 & 0.928 \\
DS-HMamba w/o entropy weighting & 0.911 & 0.913 & 0.930 \\
\textbf{Heimdall (full)}        & \textbf{0.927} & \textbf{0.929} & \textbf{0.943} \\
\bottomrule
\end{tabular}
\end{table}

Table~\ref{tab:ablation} ablates the four design choices that separate
Heimdall from a vanilla Mamba.  Each component contributes:
\begin{itemize}
\item Adding the inner scan to the outer scan (full DS-HMamba w/o gate
      versus Outer-only) lifts coarse accuracy by $+2.0$ points and
      fine by $+1.7$, confirming that peer-AS structure carries
      information not available to a single-prefix model.
\item Replacing mean-pool fusion with the AS-conditioned gate lifts
      another $+2.6$ at coarse and $+1.5$ at fine, showing that
      \emph{which} of the two scans matters depends on the AS.
\item Per-bit entropy weighting yields a final $+1.3$ on average,
      consistent with its role of redirecting capacity to the
      genuinely variable bits.
\end{itemize}

\paragraph{Sampling temperature.}  Following Gungnir we sweep
$T\in\{0.2,0.4,0.6,0.8\}$.  Best end-to-end coverage is reached at
$T=0.6$, matching Gungnir's recommendation.  However, when the
adaptive controller is enabled the closed-loop variant becomes
\emph{robust} to $T_{0}$ in $[0.4,1.0]$: starting at any value in this
band, the temperature converges to $\sim 0.55$ within five rounds.

% =======================================================================
\section{Long-Term Multi-Protocol Probing and FRP Analysis}
\label{sec:long}

We ran Heimdall continuously for three months over $23$ active
protocols (ICMPv6, TCP, UDP).  Compared against the publicly released
hitlist~\cite{ipv6hitlist} and Luori's FRP list, our crawl extends
\textsc{Cov}\textsubscript{AS} from $335$ / $530$ to $\mathbf{684}$
and \textsc{Cov}\textsubscript{Routing} from $3{,}237$ / $7{,}324$ to
$\mathbf{9{,}512}$.  Three findings echo and refine the observations
of~\cite{gungnir2025}:

\paragraph{Concentration.}  Three organisations---Cloudflare, Amazon,
and Akamai---contribute $>$2\,M ICMPv6 FRPs, dominated by
fine-granularity $/64$--$/56$ prefixes consistent with
\texttt{ip\_nonlocal\_bind} mass-binding.

\paragraph{Service vs.\ port mismatch.}  Of $1.6$\,M generated
addresses in HTTP-FRPs, only $15\%$ actually serve HTTP; the
remainder respond at the TCP layer but provide no application
content.  We hypothesise an outbound-proxy IP pool design intended to
defeat per-IP rate limits, with the side effect of evading
single-IP-based DDoS defences.

\paragraph{Security tells.}  Among SSH-FRPs we observed
\textsc{OpenSSH 7.2p2} retaining \texttt{diffie-hellman-group14-sha1}
in $\sim$11\% of fingerprints, strongly suggesting honeypot
deployments.  Among SSL/TLS-FRPs, $\sim$23\% share a single Cloudflare
wildcard certificate; private-key compromise would simultaneously
expose every covered FRP, motivating automated rotation.

% =======================================================================
\section{Ethical Considerations}\label{sec:ethics}
We follow the Menlo report~\cite{menlo} and the community guidance
of~\cite{partridge2016ethical}.  All scanning runs from a single
vantage point at $\leq 1$\,Mbps with randomised destination ordering
and no single batch exceeding 10 minutes.  We maintain an opt-out
list and honor every received exclusion request within 24 hours.  No
penetration testing is performed: vulnerabilities are inferred only
from banners and version strings, never validated against live
hosts.

% =======================================================================
\section{Conclusion}\label{sec:concl}
We presented \textsc{Heimdall}, a closed-loop FRP discovery system
that replaces Transformer-based autoregression with a \emph{Dual-Scan
Hierarchical Mamba}.  By scanning bits within a prefix \emph{and}
samples within an AS, gating their fusion on AS metadata, training
with entropy-weighted layer-wise MLM, and running an online
APD-feedback loop with bidirectional seed filtering and adaptive
temperature, Heimdall improves over the SOTA Gungnir by up to
$1.9\times$ in discovery throughput and $1.2\times$ in AS coverage at
identical probing budgets, while halving wall-clock time.  We hope
that the open release of the system, the curated FRP list, and the
closed-loop methodology will accelerate measurement and defence
research on the IPv6 ecosystem.

% =======================================================================
\bibliographystyle{IEEEtran}
\begin{thebibliography}{99}

\bibitem{durumeric2013zmap}
Z.~Durumeric, E.~Wustrow, and J.~A. Halderman, ``ZMap: Fast Internet-wide
scanning and its security applications,'' in \emph{USENIX Security}, 2013.

\bibitem{rye2023ipv6}
E.~Rye and D.~Levin, ``IPv6 hitlists at scale: Be careful what you wish
for,'' in \emph{ACM SIGCOMM}, 2023.

\bibitem{vermeulen2020diamond}
K.~Vermeulen \emph{et~al.}, ``Diamond-Miner: Comprehensive discovery of
the Internet's topology diamonds,'' in \emph{NSDI}, 2020.

\bibitem{murdock2017target}
A.~Murdock, F.~Li, P.~Bramsen, Z.~Durumeric, and V.~Paxson, ``Target
generation for Internet-wide IPv6 scanning,'' in \emph{IMC}, 2017.

\bibitem{yang2022sixforest}
T.~Yang, Z.~Cai, B.~Hou, and T.~Zhou, ``6Forest: An ensemble
learning-based approach to target generation for Internet-wide IPv6
scanning,'' in \emph{IEEE INFOCOM}, 2022.

\bibitem{hou2021sixhit}
B.~Hou, Z.~Cai, K.~Wu, J.~Su, and Y.~Xiong, ``6Hit: A reinforcement
learning-based approach to target generation for Internet-wide IPv6
scanning,'' in \emph{IEEE INFOCOM}, 2021.

\bibitem{song2022det}
G.~Song \emph{et~al.}, ``DET: Enabling efficient probing of IPv6 active
addresses,'' \emph{IEEE/ACM Trans.\ Netw.}, 2022.

\bibitem{yang2022sixgraph}
T.~Yang, B.~Hou, Z.~Cai, K.~Wu, T.~Zhou, and C.~Wang, ``6Graph: A
graph-theoretic approach to address pattern mining for Internet-wide
IPv6 scanning,'' \emph{Computer Networks}, 2022.

\bibitem{gasser2018clusters}
O.~Gasser \emph{et~al.}, ``Clusters in the expanse: Understanding and
unbiasing IPv6 hitlists,'' in \emph{IMC}, 2018.

\bibitem{zirngibl2022rusty}
J.~Zirngibl, L.~Steger, P.~Sattler, O.~Gasser, and G.~Carle, ``Rusty
clusters? Dusting an IPv6 research foundation,'' in \emph{IMC}, 2022.

\bibitem{sattler2023packed}
P.~Sattler, J.~Zirngibl, M.~Jonker, O.~Gasser, G.~Carle, and R.~Holz,
``Packed to the brim: Investigating the impact of highly responsive
prefixes on Internet-wide measurement campaigns,'' \emph{Proc.\ ACM
Networking}, vol.~1, 2023.

\bibitem{luori2024}
D.~Cheng \emph{et~al.}, ``Luori: Active probing and evaluation of
Internet-wide IPv6 fully responsive prefixes,'' in \emph{IEEE ICNP},
2024.

\bibitem{gungnir2025}
C.~Wei, Y.~Liu, L.~He, D.~Cheng, and J.~Zhou, ``Gungnir:
Autoregressive model for unified generation of IPv6 fully responsive
prefixes,'' in \emph{IEEE ICNP}, 2025.

\bibitem{gu2023mamba}
A.~Gu and T.~Dao, ``Mamba: Linear-time sequence modeling with selective
state spaces,'' arXiv:2312.00752, 2023.

\bibitem{zhu2024vimamba}
L.~Zhu \emph{et~al.}, ``Vision Mamba: Efficient visual representation
learning with bidirectional state-space model,'' in \emph{ICML}, 2024.

\bibitem{rfc4291}
R.~Hinden and S.~Deering, ``IP version 6 addressing architecture,''
RFC 4291, 2006.

\bibitem{routeviews}
University of Oregon, ``Route Views project,''
\url{https://www.routeviews.org/}, 2024.

\bibitem{caida-as2org}
CAIDA, ``AS-to-organization mapping dataset,''
\url{https://www.caida.org/catalog/datasets/as-organizations/}, 2024.

\bibitem{ipv6hitlist}
TUM, ``IPv6 hitlist service,''
\url{https://ipv6hitlist.github.io/}, 2025.

\bibitem{menlo}
E.~Kenneally and D.~Dittrich, ``The Menlo report: Ethical principles
guiding information and communication technology research,''
\emph{SSRN}, 2014.

\bibitem{partridge2016ethical}
C.~Partridge and M.~Allman, ``Ethical considerations in network
measurement papers,'' \emph{Commun.\ ACM}, 2016.

\end{thebibliography}

\end{document}
